\documentclass[a4paper,12pt]{article}

\usepackage[utf8]{inputenc}
\usepackage{geometry}
\usepackage{graphicx}
\usepackage{amsmath}
\usepackage{hyperref}

% Language and font packages
% \usepackage{ctex} % Uncomment if using Chinese characters

% Pseudocode packages - Chose algorithm2e for simplicity, removed conflict
\usepackage[ruled,vlined]{algorithm2e} 

% Fancy style packages
\usepackage{fancyhdr}
\usepackage{color}
\usepackage{xcolor}
\usepackage{tcolorbox}
\usepackage{framed}

% Math and symbol packages
\usepackage{amsfonts}
\usepackage{amssymb}
\usepackage{mathtools}

% Table and figure packages
\usepackage{booktabs}
\usepackage{multirow}
\usepackage{caption}
\usepackage{subcaption}

% Miscellaneous
\usepackage{enumitem}
\usepackage{listings}
\usepackage{tikz}
\usepackage{pgfplots}

\title{Single's Inferno}
\author{B13901165 Ivan Chang}
\date{\today}

\begin{document}

\maketitle

\section{Background}
During the welcome camp, NTUEE and NTUDFLL always work together to design the most engaging and innovative activities for freshmen. 
One of the most crucial missions is to form the ultimate “teammate pairs.” Each student has their own unique personality, preferences, and compatibility level with others. 
To avoid chaos, awkward silence, and the existential crisis of having no partner to hold a balloon with… we need a systematic and fair way to assign pairs.
The question is: how can we find a pairing that satisfies all of the students? Does such a pairing exist?

\section{Pairing Up Team Leaders} % Fixed Typo: Paring -> Pairing
\subsection{Input Format}
During the NTU welcome camp, every incoming group needs two team leaders to guide them through a series of thrilling activities. 
There are $n$ male leaders and $n$ female leaders. Some pairs are known to be compatible, 
meaning they can effectively collaborate as co-leaders.
We model this as a bipartite graph:
\begin{itemize}
    \item The set of male leaders forms one partition. $M = \{m_1, m_2, \dots, m_n\}$ % Fixed formatting
    \item The set of female leaders forms the other partition. $F = \{f_1, f_2, \dots, f_n\}$ % Fixed formatting
    \item An edge exists between a male leader and a female leader if they are compatible.
\end{itemize}

\subsection{Input}
First line: Two integers $n$ and $m$, where $n$ is the number of males and females (totally $2n$ leaders), and $m$ is the number of compatible pairs.\\
Next $m$ lines: Each line contains two integers $u$ and $v$, indicating compatible pairs.

\subsection{Solution} 
This is a classic bipartite matching problem. We can use the Maximum Flow algorithm to find the maximum matching.
We construct a flow network as follows:
1. Create a directed graph $G' = (V', E')$(requires time complexity $O(E)$), where:
\begin{itemize}
    \item Create a source node $S$ and a sink node $T$.
    \item Connect the source $S$ to each male leader $m_i$ with an edge of capacity 1.
    \item Connect each female leader $f_j$ to the sink $T$ with an edge of capacity 1.
    \item For each compatible pair $(m_i, f_j)$, create an edge from $m_i$ to $f_j$ with capacity 1.
\end{itemize}
2. Apply the Ford-Fulkerson algorithm to find the maximum flow from $S$ to $T$ in $G'$(requires time complexity $O(VE)$).\\
3. If the maximum flow equals $n$, then a perfect matching exists
and we can pair up all leaders. Otherwise, it's impossible to pair everyone up.

\section{Now...... with preferences!}
However, in reality. students have preferences for their teammates. For each students,
they will score their possible teammates from 1 to n(n is the most preferred).
But as the organizer, you want to make suer all the pairing is \textbf{stble}.
A pairing is considered stable if there are no two leaders who would both prefer each other over their current partners.
How can we find such a stable pairing?
\subsection{Input Format}
First line: An integer $n$, the number of male and female leaders (totally $2n$ leaders).\\
Next $n$ lines: Each line contains $n$ integers representing the preference list of each male leader. 
Next $n$ lines: Each line contains $n$ integers representing the preference list of each female leader.
\subsection{Solution}
We can use the \textbf{Gale-Shapley algorithm} to find a stable matching. The algorithm works as follows:
\begin{enumerate}
    \item Initialize all leaders as single.
    \item Every male leader make a invitation to the most preferred female leader who has not yet rejected him.
    \item Each female leader consider the invitation and the temporary teammates she has.
    And choose the most preferred one. reject another.
    \item Repeat steps 2 and 3 until all leaders are paired.
\end{enumerate}
Proofment: if there is an unstable pair (m, f) in the final matching, then m has proposed to f at some point (since he prefers f over his final partner).
And f must have rejected m in favor of her final partner (since she prefers her final partner over m). This contradicts the assumption that (m, f) is an unstable pair.
Therefore, All pairs are stble in the end through Gale-Shapley algorithm.
\section{Pseudo Code}
\begin{algorithm}[H]
    \SetAlgoLined
    \DontPrintSemicolon
    \caption{Gale-Shapley (men-proposing version)}
    \KwIn{Preference lists of all men and women}
    \KwOut{A stable matching $M$ between men and women}

    Initialize all men and women as \textbf{free}\;
    \ForEach{man $m$}{
        mark all women in $m$'s preference list as \textbf{not yet proposed}\;
    }

    $M \gets \emptyset$\;

    \While{there exists a free man $m$ who has not proposed to every woman}{
        Let $w$ be the highest-ranked woman on $m$'s list
        to whom $m$ has not yet proposed\;
        Mark that $m$ has now proposed to $w$\;

        \uIf{$w$ is free}{
            Match $m$ and $w$\;
            $M \gets M \cup \{(m,w)\}$\;
            Mark $m$ and $w$ as \textbf{engaged}\;
        }
        \Else{
            Let $m'$ be $w$'s current partner in $M$\;
            \eIf{$w$ prefers $m$ over $m'$}{
                Replace $(m',w)$ by $(m,w)$ in $M$\;
                Mark $m$ as \textbf{engaged} and $m'$ as \textbf{free}\;
            }{
                $m$ remains \textbf{free}\;
            }
        }
    }
    \Return $M$\;
\end{algorithm}
\subsection{time complexity}
The time complexity of the Gale-Shapley algorithm is $O(n^2)$:
\begin{itemize}
    \item In Each iteration, at most n proposals are made by all men.
    \item each man will propose to their posible teammates at most once.
    \item Therefore, the total number of proposals is at most $n^2$.
\end{itemize}

\section{Blacklist}
However, some students are \textbf{too} annoying for some others to accept them as teammates.
To solve this, We modify the Definition of stability as follows:
A pairing is considered stable if there are no two leaders who would both prefer each other over their current partners,
and being single is also stable if other people are all in his/her blacklist.
\subsection{Input Format}
First line: An integer $n$, the number of male and female leaders (totally $2n$ leaders).\\
Next $2n$ lines: Each line contains n integers representing the preference list. if a leader has blacklisted someone,
the corresponding position in the preference list will be marked as 0.
\subsection{Solution}
We can modify the Gale-Shapley algorithm to account for blacklists:
\begin{itemize}
    \item All men will only propose to women who are not in their blacklist.
    \item All women will only consider proposals from men who are not in their blacklist.
\end{itemize}
\end{document}